% reading-guide | Status: todo | Original lines: 240-315
%!TEX root = ../main.tex
% =====================================================================
\section*{Reading guide and notation}
\addcontentsline{toc}{section}{Reading guide and notation}
% =====================================================================

\begin{keyidea}[Reading guide and notation]
\small
\textbf{Summation convention.} A repeated index, once up and once down, is summed:
$\omega_iv^i=\sum_i\omega_iv^i$.\\[2pt]
\begin{tabular}{@{}p{4.1cm}p{10.2cm}@{}}
$v^i,\ \omega_i$ & components of a vector / of a covector in a chosen basis\\
$e_i,\ e^i$ & basis of $V$ / dual basis of $V^*$, $e^i(e_j)=\delta^i{}_j$\\
$\partial_i,\ \dd x^i$ & coordinate basis of vectors / dual basis of covectors, $\dd x^i(\partial_j)=\delta^i{}_j$\\
$\langle\omega,v\rangle=\omega(v)$ & contraction: the covector $\omega$ ``eats'' the vector $v$ and returns a number\\
$\alpha\wedge\beta$ & wedge (exterior) product, antisymmetric: $\dd x\wedge\dd y=-\dd y\wedge\dd x$, $\dd x\wedge\dd x=0$\\
$\dd,\ \iota_v,\ \phi\pb,\ \Lie_X$ & exterior derivative, interior product (contraction with $v$), pull-back, Lie derivative\\
$\flat,\ \sharp,\ \hodge$ & lower an index, raise an index (both need a metric), Hodge star\\
$\partial M,\ \oint$ & boundary of $M$; integral over a closed curve\\
$ab=a\cdot b+a\wedge b$ & geometric product of vectors; $a\lc A$ left contraction, $\rev A$ reverse, $\pss$ pseudoscalar, $\dual A=A\pss^{-1}$ dual\\
$\nabla$ \ \textit{versus}\ $\vd$ & plain $\nabla_X$: covariant derivative (connection). Bold $\vd=e^i\partial_i$: the geometric-algebra vector derivative.\\
$R$ \ \textit{versus}\ $\rotor$ & $R$: Riemann curvature. $\rotor$: rotor of geometric algebra.
\end{tabular}\\[3pt]
\textbf{How derivations are written.} Every non-trivial computation is split into numbered
\emph{steps}, each with its justification, and is followed by a \emph{numerical check} whenever
possible. If a step looks too fast, look for the check right below it.
\end{keyidea}

\subhead{The colour code of the boxes}
Every concept is presented in layers. The colour of a box tells you which layer you are reading:

\medskip
\begin{center}
\small
\renewcommand{\arraystretch}{1.5}
\begin{tabular}{@{}l p{11.2cm}@{}}
\fcolorbox{navy}{kiback}{\textcolor{navy}{\bfseries\ Key idea\ }} & the one statement to remember; rules of thumb; key points at the end of a section\\
\fcolorbox{formalfr}{formalbg}{\textcolor{formalfr}{\bfseries\ Formal idea\ }} & precise definition or statement\\
\fcolorbox{intuitfr}{intuitbg}{\textcolor{intuitfr}{\bfseries\ Intuition\ }} & the picture: what the object ``really is''\\
\fcolorbox{examplefr}{examplebg}{\textcolor{examplefr}{\bfseries\ Example\ }} & a fully worked, numbered example -- try to redo it by hand\\
\fcolorbox{specialfr}{specialbg}{\textcolor{specialfr}{\bfseries\ Special case\ }} & a low-dimensional or degenerate case that shows the mechanism or checks the general formula\\
\fcolorbox{remarkfr}{remarkbg}{\textcolor{remarkfr}{\bfseries\ Remark\ }} & a side comment, convention or generalisation\\
\fcolorbox{pitfr}{pitbg}{\textcolor{pitfr}{\bfseries\ Careful\ }} & a classic trap\\
\fcolorbox{synthfr}{synthbg}{\textcolor{synthfr}{\bfseries\ Synthesis\ }} & a bridge between two or more of the languages used in this document\\
\fcolorbox{exfr}{exbg}{\textcolor{exfr}{\bfseries\ Exercises\ }} & problems to test understanding
\end{tabular}
\end{center}
Theorems, propositions, lemmas and definitions are numbered together with examples, remarks and special
cases (one counter per section), so ``Example~3.2'' and ``Theorem~3.1'' never share a number.

\subhead{Provenance: how the three source texts were merged}
This document consolidates three texts: \textbf{(A)} notes on the integration of $1$-forms, contraction, Stokes and Gauss,
with physics and residues; \textbf{(B)} a companion to a course on the mathematical foundations of general relativity
(vector spaces to curvature and geodesic deviation); and \textbf{(C)} notes building geometric algebra from the tensor
product up to the exterior derivative. Their material overlaps strongly (dual space, tensor product, metric, wedge product,
contraction, Hodge star, Stokes, curvature versus $\dd^2=0$), so each topic now appears \emph{once}, in the order in which
it is logically needed, and the overlaps are discussed in the synthesis sections.

Phrases such as ``the notes'' or ``lecture notes I4'' inside the material that comes from source~(B) refer to the lecture notes of that course; they have been kept because the examples are cited by them.

\begin{center}
\small
\begin{tabular}{@{}p{4.3cm}p{6.6cm}p{3.2cm}@{}}
\toprule
\textbf{Part} & \textbf{Theme} & \textbf{Sections}\\
\midrule
I \ Algebra & vectors, duals, tensors, metric, wedge, Clifford algebra, rotors, duality & \ref{sec:vectors}--\ref{sec:synth1}\\
II \ Manifolds and forms & charts, tangent spaces, $x,\dd x,\partial_x$, tensor fields, $\dd$, $\Lie_X$, vector derivative & \ref{sec:manifolds}--\ref{sec:synth2}\\
III \ Integration & $\int$ of $1$-forms and $k$-forms, Stokes, Gauss & \ref{sec:int1}--\ref{sec:gauss}\\
IV \ Geometry of a metric & connection, Levi-Civita, geodesics, curvature, deviation & \ref{sec:conn}--\ref{sec:synth3}\\
V \ Physics & intrinsic versus metric-dependent quantities; GR bridge & \ref{sec:phys}--\ref{sec:bridge}\\
VI \ Complex analysis & residue theorem as Stokes' theorem & \ref{sec:complex}\\
VII \ Perspective & current state, master summary, references & \ref{sec:current}--\ref{sec:summary}\\
\bottomrule
\end{tabular}
\end{center}
